\chapter{Ressemblances et différences au sein de l'espèce humaine}

\noindent\textit{Tous les êtres humains appartiennent à la même espèce, mais aucun
individu n'est identique à un autre (sauf de vrais jumeaux). Comprendre pourquoi
suppose de distinguer ce qui vient de l'hérédité de ce qui vient du mode de vie et du
milieu : c'est le point de départ de toute la génétique étudiée dans ce module.}
\vspace{8pt}

\connecteBanner

\section{Les ressemblances : les caractères de l'espèce humaine}

\begin{definition}[Caractère]
Un \textbf{caractère} est un trait observable (morphologique, physiologique\dots) qui
permet de décrire un individu : couleur des yeux, groupe sanguin, taille, présence de
deux bras, respiration pulmonaire, etc.
\end{definition}

\begin{propriete}[Caractères de l'espèce]
Certains caractères sont communs à \textbf{tous} les individus de l'espèce humaine :
squelette interne, station debout possible, respiration pulmonaire, présence d'un
cerveau développé, 46 chromosomes par cellule\dots\ Ces caractères permettent de
reconnaître un individu comme appartenant à l'espèce humaine et le distinguent des
autres espèces (chimpanzé, gorille\dots).
\end{propriete}

\begin{exemple}
Un nouveau-né camerounais et un nouveau-né norvégien, bien que très différents
d'aspect, partagent tous les caractères de l'espèce humaine : ce sont deux
représentants de la même espèce, capables de se reproduire entre eux (via leurs
descendances) et de donner une descendance féconde.
\end{exemple}

\section{Les différences entre individus}

\begin{propriete}[Une immense variabilité]
Au sein de l'espèce humaine, chaque individu présente une combinaison de caractères qui
lui est propre : couleur de peau, forme du visage, groupe sanguin, taille, aptitudes
physiques\dots\ Cette \textbf{variabilité individuelle} est la règle ; seuls les vrais
jumeaux (issus d'un même \oe uf) se ressemblent presque parfaitement.
\end{propriete}

\begin{illus}
\begin{tikzpicture}[scale=0.85]
\foreach \x/\c/\lbl in {0/onbo/A, 2.4/onbo2/B, 4.8/onbgreen/C, 7.2/onbink/D}{
  \draw[\c,fill=\c,fill opacity=0.15,line width=1.1pt] (\x,0) circle (0.8);
  \node[font=\small] at (\x,0) {\lbl};
}
\node[below,font=\scriptsize,align=center] at (0,-1) {taille\\élevée};
\node[below,font=\scriptsize,align=center] at (2.4,-1) {peau\\claire};
\node[below,font=\scriptsize,align=center] at (4.8,-1) {groupe\\O};
\node[below,font=\scriptsize,align=center] at (7.2,-1) {cheveux\\bouclés};
\end{tikzpicture}
\fig{Figure — Quatre individus (A, B, C, D) de la même espèce, chacun porteur d'une
combinaison unique de caractères.}
\end{illus}

\section{Deux origines possibles pour un caractère}

\begin{definition}[Caractère héréditaire]
Un \textbf{caractère héréditaire} est transmis des parents aux enfants par
l'intermédiaire des \textbf{gènes} (voir chapitres suivants) : groupe sanguin, couleur
des yeux, certaines maladies génétiques\dots\ Il ne change pas au cours de la vie sous
l'effet du milieu.
\end{definition}

\begin{definition}[Caractère modifié par l'environnement]
Un caractère peut aussi être \textbf{modifié par le milieu de vie} : l'alimentation,
l'exposition au soleil, l'activité physique, une maladie, un accident\dots\ On parle
parfois de caractère \textbf{acquis} : bronzage, musculature développée par le sport,
cicatrice, taille influencée par la nutrition durant l'enfance.
\end{definition}

\begin{propriete}[Un caractère peut combiner les deux]
De nombreux caractères résultent d'une \textbf{combinaison} entre un potentiel
héréditaire (fixé par les gènes reçus des parents) et l'influence du milieu. Par
exemple, la taille adulte dépend à la fois de gènes hérités des parents et de
l'alimentation pendant la croissance.
\end{propriete}

\begin{exemple}
Le bronzage de la peau après une exposition au soleil est un caractère modifié par
l'environnement : il disparaît progressivement à l'ombre et n'est pas transmis à la
descendance. La couleur naturelle de la peau, elle, est un caractère héréditaire.
\end{exemple}

\begin{remarque}
Piège fréquent : un caractère héréditaire n'est pas forcément présent à la naissance
(certaines maladies génétiques ne se manifestent qu'à l'âge adulte), et un caractère
acquis peut être durable (une cicatrice). Le critère décisif n'est pas le moment
d'apparition mais la \textbf{transmissibilité à la descendance}.
\end{remarque}

% ═══════════════════════════════════════════════════════════════
\section{L'essentiel à retenir}

\begin{synthese}
\textbf{Caractère.} Trait observable décrivant un individu.\\[4pt]
\textbf{Caractères de l'espèce.} Communs à tous les individus (squelette interne,
respiration pulmonaire\dots) : ils définissent l'appartenance à l'espèce humaine.\\[4pt]
\textbf{Variabilité individuelle.} Chaque individu porte une combinaison unique de
caractères (sauf vrais jumeaux).\\[4pt]
\textbf{Caractère héréditaire.} Transmis par les gènes des parents aux enfants,
indépendant du milieu.\\[4pt]
\textbf{Caractère modifié par l'environnement (acquis).} Résulte de l'action du milieu
de vie (alimentation, soleil, activité\dots), non transmis à la descendance.\\[4pt]
\textbf{Piège fréquent.} Le critère de distinction est la transmissibilité à la
descendance, pas le moment d'apparition du caractère.
\end{synthese}

% ═══════════════════════════════════════════════════════════════
\section{Méthodes \& savoir-faire}

\methode{Reconnaître un caractère de l'espèce}
Vérifier si le caractère est présent chez \textbf{tous} les individus humains, sans
exception : dans ce cas, il définit l'espèce.
\begin{exemple}
La respiration pulmonaire est un caractère de l'espèce humaine : tous les êtres
humains respirent avec des poumons.
\end{exemple}

\methode{Distinguer caractère héréditaire et caractère acquis}
Se demander si le caractère peut être transmis à la descendance (héréditaire) ou s'il
résulte d'une action du milieu sur l'individu, non transmissible (acquis).
\begin{exemple}
La couleur des yeux est héréditaire ; une brûlure est un caractère acquis, lié à un
accident, non transmissible.
\end{exemple}

% ═══════════════════════════════════════════════════════════════
\section{Exercices d'application}
\begin{multicols}{2}

\rubrique{Caractères de l'espèce}
\exo{}\diff{1} Donner un exemple de caractère commun à tous les êtres humains.

\exo{}\diff{2} Pourquoi la respiration pulmonaire est-elle un caractère de l'espèce ?

\rubrique{Variabilité individuelle}
\exo{}\diff{1} Citer deux caractères qui varient d'un individu à l'autre.

\exo{}\diff{2} Pourquoi les vrais jumeaux se ressemblent-ils presque parfaitement ?

\rubrique{Hérédité et environnement}
\exo{}\diff{1} Le groupe sanguin est-il un caractère héréditaire ou acquis ?

\exo{}\diff{2} Le bronzage est-il transmis à la descendance ? Justifier.

\exo{}\diff{2} Citer un exemple de caractère qui combine hérédité et environnement.

% ═══════════════════════════════════════════════════════════════
\end{multicols}

\section{Exercices d'entraînement}
\begin{multicols}{2}

\rubrique{Caractères de l'espèce}
\exo{}\diff{3} Expliquer pourquoi un caractère de l'espèce ne permet pas, à lui seul,
de distinguer deux individus entre eux.

\rubrique{Variabilité individuelle}
\exo{}\diff{2} Un élève affirme que deux frères non jumeaux devraient se ressembler
parfaitement puisqu'ils ont les mêmes parents. Expliquer pourquoi c'est faux.

\rubrique{Hérédité et environnement}
\exo{}\diff{3} Un athlète développe une musculature importante grâce à l'entraînement.
Son fils naît-il avec cette musculature ? Justifier à l'aide du vocabulaire du
chapitre.

\exo{}\diff{3} La taille adulte dépend-elle uniquement des gènes hérités des parents ?
Justifier.

\exo{}\diff{2} Une maladie génétique ne se manifeste qu'à 40 ans. Est-elle pour autant
un caractère acquis ? Justifier.

\exo{}\diff{3} Proposer une expérience de pensée permettant de vérifier si un caractère
donné est héréditaire ou acquis.

% ═══════════════════════════════════════════════════════════════
\end{multicols}

\section{Sujets type BEPC}

\sujetbac[Caractères de l'espèce et variabilité]
\begin{enumerate}[label=\arabic*)]
\item Donne deux exemples de caractères communs à tous les êtres humains.
\item Pourquoi ces caractères sont-ils qualifiés de « caractères de l'espèce » ?
\item Donne deux exemples de caractères qui varient d'un individu à l'autre.
\item Pourquoi les vrais jumeaux constituent-ils une exception à cette variabilité ?
\end{enumerate}

\sujetbac[Hérédité ou environnement ?]
\begin{enumerate}[label=\arabic*)]
\item Définis un caractère héréditaire.
\item Définis un caractère modifié par l'environnement.
\item Classe les caractères suivants en héréditaire ou acquis : groupe sanguin,
bronzage, cicatrice, couleur naturelle des cheveux.
\item Explique pourquoi la taille adulte peut dépendre à la fois de l'hérédité et de
l'environnement.
\end{enumerate}

% ═══════════════════════════════════════════════════════════════
\section{Corrigés}
\begin{multicols}{2}

\corrige{de l'exercice 1}
Par exemple : la respiration pulmonaire, ou la présence d'un squelette interne.

\corrige{de l'exercice 2}
Parce que tous les individus humains, sans exception, respirent au moyen de poumons.

\corrige{de l'exercice 3}
Par exemple : la couleur des yeux et le groupe sanguin.

\corrige{de l'exercice 4}
Parce qu'ils sont issus d'un même \oe uf fécondé et portent donc une information
génétique quasiment identique.

\corrige{de l'exercice 5}
C'est un caractère héréditaire, transmis par les gènes des parents.

\corrige{de l'exercice 6}
Non, il disparaît à l'ombre et n'est pas transmis : c'est un caractère modifié par
l'environnement.

\corrige{de l'exercice 7}
Par exemple : la taille adulte, qui dépend des gènes hérités et de l'alimentation.

\corrige{de l'exercice 8}
Parce qu'un caractère de l'espèce est partagé par tous les individus : il ne permet
donc pas de différencier deux personnes entre elles, contrairement aux caractères qui
varient d'un individu à l'autre.

\corrige{de l'exercice 9}
C'est faux : hormis les vrais jumeaux, chaque enfant reçoit une combinaison différente
de gènes de ses parents (brassage génétique), ce qui explique que des frères non
jumeaux ne se ressemblent pas parfaitement.

\corrige{de l'exercice 10}
Non : la musculature développée par l'entraînement est un caractère acquis, lié à
l'activité de l'athlète lui-même ; elle n'est pas transmise à sa descendance.

\corrige{de l'exercice 11}
Non, elle dépend à la fois de gènes hérités des parents et de facteurs du milieu, en
particulier l'alimentation durant la croissance.

\corrige{de l'exercice 12}
Non : le critère n'est pas le moment d'apparition mais la transmissibilité à la
descendance ; une maladie génétique tardive reste héréditaire si elle est due à un gène
transmis par les parents.

\corrige{de l'exercice 13}
Par exemple : comparer le caractère chez des vrais jumeaux élevés dans des milieux
différents ; s'il reste identique malgré des milieux différents, il est probablement
héréditaire.

\corrige{du sujet 1}
1) Par exemple : la respiration pulmonaire, le squelette interne.\\
2) Parce qu'ils sont présents chez tous les individus humains, sans exception.\\
3) Par exemple : la couleur des yeux, le groupe sanguin.\\
4) Parce qu'étant issus d'un même \oe uf, ils portent une information génétique quasiment
identique et se ressemblent donc fortement.

\corrige{du sujet 2}
1) Un caractère héréditaire est transmis des parents aux enfants par les gènes.\\
2) Un caractère modifié par l'environnement résulte de l'action du milieu de vie sur
l'individu, sans être transmis à la descendance.\\
3) Héréditaires : groupe sanguin, couleur naturelle des cheveux. Acquis : bronzage,
cicatrice.\\
4) Parce qu'elle dépend à la fois de gènes hérités des parents (potentiel de
croissance) et de facteurs du milieu comme l'alimentation pendant l'enfance.

\end{multicols}
\exercicesBanner
